\section{从 Vision 到 Representation：研究轨迹的主线}

访谈中最重要的个人线索，是“世界总不让我做 Vision，但我偏要做”。这句话背后有两层含义。第一层是求学和职业选择上的阻力：申请博士时几乎没有理想中的 vision offer，后来靠给 Tu Zhuowen 发邮件才进入 UCSD。第二层是时代叙事的阻力：当 LLM 成为中心后，computer vision 看似退到边缘。但谢赛宁的判断是，如果把 vision 理解成一种 perspective，而不是一组任务，那么它恰恰是通往真实智能的核心。

\begin{figure}[H]
\centering
\includegraphics[width=0.96\textwidth]{trajectory-timeline.png}
\caption{谢赛宁访谈中的研究轨迹：从 SJTU、UCSD、FAIR、DeepMind/NYU 到 AMI Labs。}
\end{figure}

\begin{knowledgebox}{读图：每一次选择都在筛选研究条件}
这张时间线不是履历表，而是选择逻辑图。SJTU 阶段提供通识和计算机入口；UCSD 阶段把问题聚焦到 vision 与 representation；FAIR 阶段通过 He Kaiming、ResNeXt 等经历学习顶尖研究品味；DeepMind 和 NYU 阶段让他看到视频、具身任务与组织资源的张力；AMI Labs 则试图把世界模型问题放到一个新的组织形态里推进。
\end{knowledgebox}

\subsection{五次实习与非线性研究}

博士阶段他做过 NEC Lab、Adobe、Meta、Google Research 和 DeepMind 等实习。有的实习产出了论文，有的没有明显成果。访谈里他并不把“没有产出”讲成失败，而是讲成探索：去不同组织看不同问题，理解不同研究文化，确认自己不想做什么。

这对年轻研究者尤其重要。研究路线并不是把每一步都优化成论文数量，而是让若干看似分散的经历最终形成一个内在问题。谢赛宁后来能把 vision、video、representation、world model 和 robotics 连接起来，正是因为早期没有只在一个狭窄路径上滚动优化。

\begin{importantbox}{研究不是 point estimate，而是时间积分}
访谈中他用“不要在乎每一个点上的估计”来描述研究评价。单篇论文、单次拒稿、单次实习是否成功，都只是时间轴上的点；一个研究者真正的质量，要看长期积累后的积分。
\end{importantbox}

\subsection{FAIR、He Kaiming 与 ResNeXt}

在 FAIR 的经历是整期访谈的一个关键节点。谢赛宁讲到 He Kaiming 加入 FAIR 后，自己在最后一个月实习中与他合作 ImageNet challenge，并发展出 ResNeXt 相关工作。这里的重点不是“遇到贵人”的故事，而是顶尖研究者如何把普通 idea 打磨成可扩展的 representation 方案。

ResNeXt 的思想可以简化为：在 ResNet 的基础上引入 cardinality，把一个分支扩展成多个并行 group，在相近计算量下获得更好的表征能力。谢赛宁在访谈里还把它和今天的 MoE 直觉相连：稀疏化、分组、可扩展能力，并不是今天才出现的想法。

\begin{knowledgebox}{术语消化：ResNeXt、cardinality 与 MoE 直觉}
\begin{tabularx}{\textwidth}{p{0.22\textwidth}p{0.32\textwidth}X}
\toprule
术语 & 解决的问题 & 与本期主线的关系 \\
\midrule
ResNet & 通过残差连接让深层网络更容易训练 & 代表早期视觉表征学习的核心范式。 \\
ResNeXt & 用多个并行分组分支提升网络表达能力 & 体现“同等计算量下更可扩展的 representation”。 \\
Cardinality & 并行分支或 group 的数量 & 比单纯加深/加宽更接近结构化扩展。 \\
MoE & Mixture of Experts，稀疏激活不同专家 & 与分组/稀疏扩展有相通的工程直觉。 \\
\bottomrule
\end{tabularx}
\end{knowledgebox}

\subsection{本章小结}

谢赛宁的研究轨迹可以概括为：从具体视觉任务出发，逐渐把 vision 理解成连续高维信号、层次化表征、空间/时间认知和预测性智能的总和。这个定义为后面的世界模型铺路。

